\documentclass[10pt]{article}
\usepackage[letterpaper,margin=0.83in]{geometry}
\usepackage{amsmath,amssymb,amsthm,booktabs}
\usepackage[hidelinks]{hyperref}
\newcommand{\F}{\mathbb F}
\newcommand{\Z}{\mathbb Z}
\newcommand{\Q}{\mathbb Q}
\newcommand{\sym}[2]{\left(\frac{#1}{#2}\right)_5}
\setlength{\parindent}{0pt}
\setlength{\parskip}{5pt}
\title{Three fifth powers: why 8, 9, and 42 are missing}
\author{Fifth42 project}
\date{27 November 2026 \quad\textbar\quad Short note for review}
\begin{document}
\maketitle
\vspace{-1.7em}
\textbf{Result.} Each of 8, 9, and 42 is a sum of three fifth powers modulo
every positive integer, but none is a sum of three signed integer fifth
powers. There is no bound on the size of the coordinates in the proof.
The case 8 answers a question posed by S\'ebastien Palcoux in 2015~\cite{palcoux}.

\section*{1. A pair of coordinates gives a test}
Let $\zeta$ be a primitive fifth root of unity in $K=\Q(\zeta)$ and
$\lambda=\zeta-1$. For a prime ideal $P\nmid 5$ and a unit $a$ modulo $P$,
the quintic residue symbol $\sym{a}{P}$ is the fifth root of unity whose
reduction is $a^{(NP-1)/5}$, where $NP$ is the residue field size.
Multiply the symbols with their ideal multiplicities for a composite
denominator. An element in the denominator denotes its principal ideal.

Suppose $u^5+v^5+w^5=m$ for $m\in\{8,9,42\}$ and $5\nmid u+v$.
Put $L=u+\zeta v$ and $h=v/(u+v)$ in $\F_5$. Then
\begin{equation}
  \sym{L}{m}=\sym{\zeta}{m}^{h}. \tag{1}
\end{equation}
Indeed, $L$ is prime to $5$ since it reduces to $u+v$ at the unique prime
above $5$. Every prime $p$ dividing $m$ is among $2,3,7$, which have order
four modulo $5$. Thus $p$ is inert in $K$, and the reduction of $\zeta$ is
not in $\F_p$. If $L$ vanished modulo the prime over $p$, both $u$ and $v$
would be divisible by $p$, and then so would $w$. That would give $p^5\mid m$,
which is impossible. At any prime ideal $P$ dividing $L$, the pair cancels
in fifth powers and $m\equiv w^5\not\equiv0\pmod P$. Thus $\sym{m}{L}=1$.

Choose an integer representative for $h$. The element
$\gamma=\zeta^{-h}L$ satisfies
$\gamma\equiv u+v+\lambda(v-h(u+v))\equiv u+v\pmod{\lambda^2}$,
so it is primary. Eisenstein reciprocity~\cite[Eq.~(1.1)]{sw} gives
$\sym{\gamma}{m}=\sym{m}{\gamma}=\sym{m}{L}=1$, proving (1).
If $\gamma$ is a unit, first apply the nonunit formulation to $q^5\gamma$
for any rational prime $q\nmid 5m$; the extra factor contributes $1$ to
both symbols.

\section*{2. Start with 8}
The fifth powers of $0,1,2,3,4$ modulo 25 are $0,1,7,-7,-1$,
respectively; the base modulo 5 determines the fifth power modulo 25.
The only unordered sum congruent to 8 takes $0,1,7$. Choose the pair whose
bases are $2,1$. Then $h=1/3=2$ in $\F_5$.

For $p=2,3,7$, write $T$ for the reduction of $\zeta$ in
$\F_{p^4}=\F_p[T]/(1+T+T^2+T^3+T^4)$. Denote by $e_p(u,v)$ the exponent
of $T$ in $(u+Tv)^{(p^4-1)/5}$. For $m=8=2^3$, (1) gives
$e_2(u,v)=2e_2(0,1)=2\cdot3=1$ modulo 5, after cancelling the factor 3.
But the pairs $(1,0),(0,1),(1,1)$ modulo 2 have exponents $0,3,4$.
In particular $(1+T)^3=1+T+T^2+T^3=T^4$. So 8 has no representation.

\clearpage
\section*{3. The same test for 9 and 42}
\textbf{9.} Modulo 25 the only unordered sum for 9 uses $1,1,7$. Choose
the pair with bases $1,2$, so $h=2/3=4$ in $\F_5$. For $m=9=3^2$,
(1) requires $e_3(u,v)=4e_3(0,1)=4$, after cancelling the factor 2.
Scalar multiplication by a nonzero element of $\F_3$ leaves the exponent
unchanged, since such scalars are fifth powers. The four pairs up to
scaling, $(1,0),(0,1),(1,1),(1,-1)$, have exponents $0,1,3,3$.
For the last two values, $T^9=T^{-1}$ gives
$(1\pm T)^8=\pm T^{-1}$ and hence $(1\pm T)^{16}=T^3$.
The required exponent 4 is absent.

\textbf{42.} Modulo 25 the only unordered sum for 42 uses $-7,-1,0$,
so order the bases as $(A,B,C)\equiv(3,4,0)\pmod5$.
Define $S=e_2+e_3+e_7$. The exponents $e_p(0,1)$ are $3,1,0$,
respectively. Equation (1) says $S(u,v)=4h$ modulo 5. The pairs
$(A,C),(B,C),(A,B)$ have $h=0,0,2$, respectively, so
\[
  -S(A,C)+S(B,C)+S(A,B)\equiv3\pmod5.
\]
At each prime $p\mid42$, the triple is nonzero and its fifth powers sum
to zero modulo $p$. Let $E_p=-e_p(A,C)+e_p(B,C)+e_p(A,B)$.
Here is the complete finite calculation, with exponents taken modulo 5:
\begin{center}
\begin{tabular}{@{}cclc@{}}
\toprule
$p$ & $e_p(0,1)$ & $e_p(1,r)$ for $r=0,\ldots,p-1$ & Possible $E_p$\\
\midrule
2 & 3 & $0,4$ & $4$\\
3 & 1 & $0,3,3$ & $3$\\
7 & 0 & $0,0,1,2,4,3,0$ & $0,2,3$\\
\bottomrule
\end{tabular}
\end{center}
The middle column of lists comes from raising $1+rT$ to $3,16,480$,
respectively, with $1+T+T^2+T^3+T^4=0$. For the last column, scale
the first nonzero entry to 1. Fifth powers permute $\F_p$, so the
normalized triples are $(0,1,-1)$ and the $p$ triples $(1,b,c)$ with
$c^5=-1-b^5$. For nonzero $u$, use $e_p(u,v)=e_p(1,v/u)$.
This yields only $4+3+\{0,2,3\}=\{0,2,4\}$ modulo 5, never 3.

\section*{4. Why no ordinary modulus can rule them out}
Kowalczyk and Miska~\cite[Table~3]{km} show that at every prime other
than 11, a sum of fifth powers in $\Z_p$ can be written with at most
three summands. Each target is such a sum (use copies of $1^5$ and pad
with zeros when needed). At 11, fifth powers are $0,\pm1$ modulo 11.
The targets reduce to $-3,-2,-2$, respectively, and have witnesses
$(-1,-1,-1)$ and $(-1,-1,0)$. Each lifts by Hensel's lemma since the
derivative in a unit coordinate is nonzero modulo 11. Reduce the
$\Z_p$ solutions modulo prime powers and combine them by the Chinese
remainder theorem.

\small
\textit{Scope.} This AI-assisted note is for mathematical review. It treats
integer coordinates only. The \href{https://github.com/paulnord/fifth42/blob/494c8a3a85ecce84505fd47aabadc9465303393c/docs/fifth42.pdf}{full manuscript}
contains the rational-denominator extension and verification materials.

\begin{thebibliography}{9}
\setlength{\itemsep}{0pt}
\bibitem{palcoux} S. Palcoux, \href{https://math.stackexchange.com/questions/1524797/can-a-sum-of-three-fifth-power-of-integers-be-8}{Can a sum of three fifth power of integers be 8?}, Mathematics Stack Exchange, 11 November 2015.
\bibitem{sw} B. K. Spearman and K. S. Williams, \href{https://people.math.carleton.ca/~williams/papers/pdf/142.pdf}{A simple proof of Eisenstein's reciprocity law from Stickelberger's theorem}, Indian J. Pure Appl. Math. \textbf{17} (1986), 169--174.
\bibitem{km} T. Kowalczyk and P. Miska, \href{https://doi.org/10.1112/mtk.12276}{On Waring numbers of henselian rings}, Mathematika \textbf{70} (2024), e12276.
\end{thebibliography}
\end{document}
